\section{Proof details and analytic firewalls}
\label{app:proofs}

\subsection{Why the row decomposition is nuclear rather than entrywise}

An entrywise \(\ell^1\) estimate would sum

\[
 \sum_{d\ge2}\sum_{q\ge q_0(d)}
 d^{-\sigma}(dq-1)^{-\sigma}q^{-a},
\]

which imposes unnecessarily strong exponents.  The correct structure is that
all entries with the same output \(d\) form one rank-one row.  The nuclear
cost of the row is its \(\ell^2\), not \(\ell^1\), norm.  This is the source
of the exact inner threshold \(2(\sigma+a)>1\), followed by the outer
threshold \(2\sigma>1\).

For uniform constants, \(dq\ge3\) gives

\[
 (dq)^{-2\sigma}
 \le (dq-1)^{-2\sigma}
 \le (3/2)^{2|\sigma|}(dq)^{-2\sigma}
\]

after choosing the inequality direction according to the sign and absorbing
it into a compact-parameter constant.  Thus on every compact subset of the
claimed domain,

\[
 C_1d^{-4\sigma}
 \le\|R_d\|_1^2
 \le C_2d^{-4\sigma}.
\]

The \(q=1\) lower bound applies for every \(d\ge3\); the exceptional row
\(d=2\) does not affect outer summability.

\subsection{Trace-norm holomorphy}

Let \(K\) be a compact subset of

\[
 \Re s>\frac12,
 \qquad
 \Re(s+u)>\frac12.
\]

There is an \(\varepsilon>0\) such that both real parts exceed their
boundaries by \(2\varepsilon\) throughout \(K\).  A derivative of order
\((j,k)\) in \((s,u)\) multiplies an entry by a polynomial in

\[
 \log d,\quad\log(dq-1),\quad\log q.
\]

For every fixed polynomial degree these logarithms are bounded by a constant
times \((dq)^\varepsilon\).  Replacing the margin \(2\varepsilon\) by
\(\varepsilon\) reproduces a summable row majorant.  The differentiated row
series therefore converges locally uniformly in trace norm.  This proves
Fréchet holomorphy into \(\Sone\), and the standard determinant map is
holomorphic on that ideal.

\subsection{The Fourier projection on ideals}

The diagonal unitary action

\[
 U_te_n=\e^{int}e_n
\]

acts isometrically on every Schatten ideal by conjugation.  Bochner
integration therefore gives a contraction

\[
 \mathcal P_j(T)=\frac1{2\pi}\int_0^{2\pi}
 \e^{-ijt}U_tTU_t^{-1}\,\dd t.
\]

It retains the matrix diagonal \(d-n=j\).  For \(j=1\), the graph relation
forces \(d=n+1\), whose cofactor is one.  This proves both the scalar
necessity and the pure-cofactor noncompactness.  The same averaging preserves
the compact ideal because it is an operator-norm closed convex ideal.

\subsection{A bounded pure-cofactor regime}

Put \(a=\Re u>1\).  For fixed output \(d\), sources are \(dq-1\), so the
absolute row sum is at most

\[
 \sum_{q\ge1}q^{-a}=\zeta(a).
\]

For fixed input \(n\), every allowed edge corresponds to a divisor
\(q\mid n+1\), and the absolute column sum is at most the same full series.
Schur's test gives \(\|Q_u\|\le\zeta(a)\).  Combined with
\cref{thm:pure-obstruction}, this provides an explicit nonempty regime in
which the matrix is bounded but noncompact.  No claim is made for the entire
strip \(1/2<a\le1\).

\subsection{Character coefficients and local logarithms}

In the trace-class domain, \(D_\chi(s,z)\) is entire in \(z\), but its
logarithm need not be globally single-valued across zeros.  For
\(|z|\|L_{s,\chi}\|<1\), the normalized logarithm is uniquely determined by

\[
 -\log D_\chi(s,z)=
 \sum_{r\ge1}\frac{z^r}{r}\Tr L_{s,\chi}^r.
\]

Every fixed-period trace has finite group support by
\cref{lem:confinement}.  Haar coefficient extraction may therefore be
performed term by term.  This local procedure is all that is needed for
\cref{thm:coefficient}; no global logarithm branch or primitive-product
convergence at \(z=1\) is used.

\subsection{Radius of the isolated holonomy-two series}

Stirling's formula gives

\[
 \log M_k
 =\log\frac{(2k-1)!}{(k-1)!}
 =k\log k+O(k).
\]

Hence for \(\sigma=\Re s\),

\[
 \left|M_k^{-2s}\right|^{1/k}
 =\exp[-2\sigma\log k+O(1)].
\]

The root tends to zero for \(\sigma>0\), to one for \(\sigma=0\), and to
infinity for \(\sigma<0\).  This proves the stated radii.  The calculation
belongs to one Fourier coefficient and does not enlarge the whole-operator
domain \cref{eq:S1-domain}.

\subsection{Regular lift in the semifinite algebra}

The row estimate for \(\mathbb R_d\) proves more than formal integrability:
the sum of row operator norms is bounded by the same summable row
\(\ell^2\)-norms for \(\Re s>1/2\).  Thus the lift is a bounded element of
\(\cN\) as well as an element of \(L^1(\cN,\Phi)\) there.  Its powers remain
in \(L^1\), and the local trace series in \cref{eq:phi-det} is meaningful.

Ordinary noncompactness is independent of this estimate.  If the chosen
vector in \cref{thm:lift-noncompact} is first taken with finite deck support,
an escaping sequence of right translates is eventually orthogonal.  A
general nonzero image is approximated by such a vector.  This gives the
required weakly null witness without assuming a spectral decomposition of
the lift.

\subsection{First-return and repetition conventions}

The induced return section consists of edges with \(q>1\).  Each \(C_k\)
visits it once, so inducing turns a length-\(k\) primitive orbit into one fixed
branch with a retained time marker \(z^k\).  Setting \(z=1\) forgets that
time; it does not make the branches identical as source histories, but it
does make the pure cofactor diagonal entries identical.  Endpoint
regularization retains the full source history through \(M_k^{-2s}\).

This is distinct from a temporal repetition.  A repetition of an orbit
raises its holonomy to a power.  Atomic holonomy therefore admits no
repetitions by \cref{cor:no-atomic-repeat}, even though it has infinitely many
different primitive representatives indexed by \(k\).
